\begin{longtable}{p{0.58\textwidth}p{0.34\textwidth}}
\caption*{源程序与说明}\\
\toprule 文件 & 用途\\
\midrule\endfirsthead
\toprule 文件 & 用途\\
\midrule\endhead
AI工具使用详情.pdf & AI工具使用过程、采纳与人工核验说明\\
decompose\_effects.py & 模型求解、验证、出图或材料整理源程序\\
export\_results.py & 模型求解、验证、出图或材料整理源程序\\
extend\_long\_results.py & 模型求解、验证、出图或材料整理源程序\\
generate\_paper\_data.py & 模型求解、验证、出图或材料整理源程序\\
model\_core.py & 模型求解、验证、出图或材料整理源程序\\
plot\_results.py & 模型求解、验证、出图或材料整理源程序\\
prepare\_support.py & 模型求解、验证、出图或材料整理源程序\\
robustness\_analysis.py & 模型求解、验证、出图或材料整理源程序\\
sensitivity\_analysis.py & 模型求解、验证、出图或材料整理源程序\\
solve\_all.py & 模型求解、验证、出图或材料整理源程序\\
validate\_model.py & 模型求解、验证、出图或材料整理源程序\\
verify\_results.py & 模型求解、验证、出图或材料整理源程序\\
建模结果汇总.md & 建模结果与复核摘要\\
\bottomrule
\end{longtable}

\begin{longtable}{p{0.58\textwidth}p{0.34\textwidth}}
\caption*{输入数据}\\
\toprule 文件 & 用途\\
\midrule\endfirsthead
\toprule 文件 & 用途\\
\midrule\endhead
附件/附件1.xlsx & 模型运行所需的题目输入数据\\
附件/附件2.xlsx & 模型运行所需的题目输入数据\\
\bottomrule
\end{longtable}

\begin{longtable}{p{0.58\textwidth}p{0.34\textwidth}}
\caption*{图表文件}\\
\toprule 文件 & 用途\\
\midrule\endfirsthead
\toprule 文件 & 用途\\
\midrule\endhead
figures/fig01\_烘房环境边界.pdf & 论文图表的位图或矢量版本\\
figures/fig01\_烘房环境边界.png & 论文图表的位图或矢量版本\\
figures/fig02\_半径收缩曲线.pdf & 论文图表的位图或矢量版本\\
figures/fig02\_半径收缩曲线.png & 论文图表的位图或矢量版本\\
figures/fig03\_问题1温度云图.pdf & 论文图表的位图或矢量版本\\
figures/fig03\_问题1温度云图.png & 论文图表的位图或矢量版本\\
figures/fig04\_问题1水分云图.pdf & 论文图表的位图或矢量版本\\
figures/fig04\_问题1水分云图.png & 论文图表的位图或矢量版本\\
figures/fig05\_问题1径向剖面.pdf & 论文图表的位图或矢量版本\\
figures/fig05\_问题1径向剖面.png & 论文图表的位图或矢量版本\\
figures/fig06\_问题2温度云图.pdf & 论文图表的位图或矢量版本\\
figures/fig06\_问题2温度云图.png & 论文图表的位图或矢量版本\\
figures/fig07\_问题2水分云图.pdf & 论文图表的位图或矢量版本\\
figures/fig07\_问题2水分云图.png & 论文图表的位图或矢量版本\\
figures/fig08\_问题2径向剖面.pdf & 论文图表的位图或矢量版本\\
figures/fig08\_问题2径向剖面.png & 论文图表的位图或矢量版本\\
figures/fig09\_问题3水分剖面.pdf & 论文图表的位图或矢量版本\\
figures/fig09\_问题3水分剖面.png & 论文图表的位图或矢量版本\\
figures/fig10\_问题3水分历程.pdf & 论文图表的位图或矢量版本\\
figures/fig10\_问题3水分历程.png & 论文图表的位图或矢量版本\\
figures/fig11\_问题4水分剖面.pdf & 论文图表的位图或矢量版本\\
figures/fig11\_问题4水分剖面.png & 论文图表的位图或矢量版本\\
figures/fig12\_问题4水分历程.pdf & 论文图表的位图或矢量版本\\
figures/fig12\_问题4水分历程.png & 论文图表的位图或矢量版本\\
figures/fig13\_固定与收缩模型比较.pdf & 论文图表的位图或矢量版本\\
figures/fig13\_固定与收缩模型比较.png & 论文图表的位图或矢量版本\\
figures/fig14\_网格收敛.pdf & 论文图表的位图或矢量版本\\
figures/fig14\_网格收敛.png & 论文图表的位图或矢量版本\\
figures/fig15\_参数灵敏度.pdf & 论文图表的位图或矢量版本\\
figures/fig15\_参数灵敏度.png & 论文图表的位图或矢量版本\\
figures/fig16\_边界口径稳健性.pdf & 论文图表的位图或矢量版本\\
figures/fig16\_边界口径稳健性.png & 论文图表的位图或矢量版本\\
figures/fig17\_扩散系数机理.pdf & 论文图表的位图或矢量版本\\
figures/fig17\_扩散系数机理.png & 论文图表的位图或矢量版本\\
\bottomrule
\end{longtable}

\begin{longtable}{p{0.58\textwidth}p{0.34\textwidth}}
\caption*{结果与核验文件}\\
\toprule 文件 & 用途\\
\midrule\endfirsthead
\toprule 文件 & 用途\\
\midrule\endhead
results/mechanism\_decomposition.json & 正式结果表、压缩场快照或校验记录\\
results/q1\_paper\_snapshot.npz & 正式结果表、压缩场快照或校验记录\\
results/q2\_paper\_snapshot.npz & 正式结果表、压缩场快照或校验记录\\
results/q3\_paper\_snapshot.npz & 正式结果表、压缩场快照或校验记录\\
results/q4\_paper\_snapshot.npz & 正式结果表、压缩场快照或校验记录\\
results/result1.xlsx & 正式结果表、压缩场快照或校验记录\\
results/result2.xlsx & 正式结果表、压缩场快照或校验记录\\
results/result3.xlsx & 正式结果表、压缩场快照或校验记录\\
results/result4.xlsx & 正式结果表、压缩场快照或校验记录\\
results/robustness.csv & 正式结果表、压缩场快照或校验记录\\
results/robustness.json & 正式结果表、压缩场快照或校验记录\\
results/sensitivity.csv & 正式结果表、压缩场快照或校验记录\\
results/sensitivity.json & 正式结果表、压缩场快照或校验记录\\
results/validation\_summary.json & 正式结果表、压缩场快照或校验记录\\
results/xlsx\_validation.json & 正式结果表、压缩场快照或校验记录\\
\bottomrule
\end{longtable}
